\chapter{Contexte et État de l'art}

\section{Contexte du projet}

Le football, fortement transformé par la digitalisation du sport, produit aujourd’hui une grande quantité de données liées aux équipes, aux joueurs et aux matchs. Malgré cette richesse, leur exploitation reste complexe en raison du caractère imprévisible du jeu et de la multiplicité des facteurs influençant les performances, tels que la forme des équipes, les statistiques ou les conditions de jeu. Cette complexité limite les approches classiques et rend pertinente l’utilisation de méthodes basées sur le machine learning.

Dans ce contexte, le projet Football Prediction Dashboard vise à exploiter les données historiques du football afin de construire une plateforme capable de prédire les résultats des matchs. Le système modélise les probabilités des différentes issues possibles et propose une visualisation claire des prédictions via une interface web interactive, permettant une meilleure interprétation des performances et tendances du jeu.

\section{État de l'art}

La prédiction des résultats de football a évolué de simples modèles statistiques vers des architectures complexes de Machine Learning. Les premiers travaux significatifs reposaient sur des distributions de Poisson pour modéliser le nombre de buts marqués par chaque équipe. Maher \cite{maher1982} a jeté les bases en utilisant des variables indépendantes pour l'attaque et la défense, tandis que Dixon et Coles \cite{dixon1997} ont amélioré ces modèles en introduisant des facteurs de pondération temporelle pour accorder plus d'importance aux performances récentes.

Avec l'avènement du Big Data, les chercheurs se sont tournés vers l'apprentissage automatique. Bunker et Thabtah \cite{bunker2019} ont démontré qu'un cadre structuré de Machine Learning surpasse les méthodes statistiques traditionnelles en capturant des relations non linéaires entre les variables. Aujourd'hui, des algorithmes comme XGBoost \cite{chen2016xgboost} sont largement adoptés pour leur capacité à traiter des volumes massifs de données tout en offrant une précision élevée grâce à l'optimisation par gradient.

Le tableau~\ref{tab:comparatif_modeles} synthétise les principales approches identifiées dans la littérature et leurs caractéristiques respectives.


\begin{table}[H]
    \centering
    \renewcommand{\arraystretch}{1.5}
    \begin{tabularx}{\textwidth}{|L|c|Y|Y|c|}
        \hline
        \rowcolor{MediumGray}
        \textbf{Modèle} & \textbf{Int.} & \textbf{Avantages} & \textbf{Limites} & \textbf{Réf.} \\
        \hline
        Régression Logistique & Haute & Rapide, simple, probabilités claires. & Difficulté avec les relations complexes. & \cite{bunker2019} \\
        \hline
        SVM & Faible & Efficace pour les données complexes. & Sensible à l'échelle des données. & \cite{cortes1995svm} \\
        \hline
        Random Forest & Moyenne & Modèle robuste, réduit le surapprentissage. & Aspect "boîte noire", lourd en mémoire. & \cite{breiman2001rf} \\
        \hline
        \rowcolor{LightGray}
        \textbf{XGBoost} & \textbf{Moyenne} & \textbf{Correction des erreurs, très haute précision.} & \textbf{Risque de surapprentissage si mal réglé.} & \cite{chen2016xgboost} \\
        \hline
    \end{tabularx}
    \caption{Tableau comparatif des modèles d'apprentissage automatique pour la prédiction sportive}
    \label{tab:comparatif_modeles}
\end{table}


\section{Objectifs du projet}

Le projet \textit{Football Prediction Dashboard} vise à développer une plateforme intelligente capable d’analyser des données historiques de football afin de prédire les résultats des matchs. Les objectifs principaux sont les suivants :

1. \textbf{Collecte et structuration des données :} Mettre en place un processus d’ingestion et de préparation de données issues de compétitions de football, afin de construire une base exploitable pour l’analyse.

2. \textbf{Modélisation prédictive :} Concevoir et exploiter des modèles de machine learning capables d’estimer les probabilités des issues possibles d’un match (victoire domicile, match nul, victoire extérieur) avec une bonne fiabilité.

3. \textbf{Visualisation des résultats :} Développer une interface web interactive permettant de présenter les prédictions de manière claire, intuitive et exploitable pour l’utilisateur.

4. \textbf{Fiabilité et exploitation des résultats :} Assurer une cohérence des prédictions en s’appuyant sur des données pertinentes et des indicateurs représentatifs des performances des équipes et des joueurs.

\section{Conclusion}

Ce chapitre a présenté le contexte du projet et les enjeux liés à l’analyse des données dans le football ainsi que la complexité de la prédiction des résultats sportifs. Il a introduit les principales approches de modélisation et leurs caractéristiques. Dans ce cadre, le projet vise à exploiter ces méthodes pour produire des prédictions fiables. Les chapitres suivants détailleront la méthodologie, la conception du système et les résultats obtenus.
